%----------------------------------------------------------------------------
\chapter{Discussion}
\label{chap:discussion}
%----------------------------------------------------------------------------

This chapter interprets the results of \autoref{chap:results} and \autoref{chap:ablations} and explains the design choices described in \autoref{sec:scope}.

%----------------------------------------------------------------------------
\section{What the soft prompt carries}
\label{sec:disc-carries}
%----------------------------------------------------------------------------

The results show that target-specific information reaches the language model through the soft prompt. The prefix is the only input that differs between papers, while the LoRA parameters are shared across papers and fixed during generation (\autoref{sec:reader}). With the same separately trained reader, the zero input produced a single text for all 945 evaluation papers, whereas the GNN estimate produced 939 distinct abstracts (\autoref{tab:ablation-results}). The identification results also indicate that the generated texts contain vocabulary specific to their target papers: compared with all 945 original abstracts, an abstract generated by the final model is most similar to its own original for 18.8\% of the papers, against 0.11\% for a random ranking (\autoref{tab:final-model}).

The soft prompt also conveys information about the target paper's research content. For each content aspect assessed separately, approximately 43--45\% of the final model's abstracts receive a score of at least 2, indicating agreement with the original abstract, possibly with differences in detail. This shows that neighbourhood-based conditioning can support descriptions that reflect specific aspects of the target's research content.

The examples in \autoref{app:examples} show that a generated abstract can capture the target's research area while changing what the paper actually does or contributes. A possible explanation is that references constrain the research context more strongly than the particular contribution: they identify related work, but do not uniquely specify how the target paper combines, modifies or evaluates it. These failures may also arise from information lost during neighbourhood encoding or abstract generation. The present results do not separate these explanations.

%----------------------------------------------------------------------------
\section{Representation quality and content agreement}
\label{sec:disc-representation}
%----------------------------------------------------------------------------

Learning to combine and transform the reference embeddings improves the content of the generated abstracts. The reference mean and the GNN estimate are computed from the same references and read by the same reader, so their comparison isolates the effect of learned estimation. With the reader held fixed, the GNN estimate raises the judge's scores for subfield, problem and task by 0.14--0.16 compared with the reference mean (\autoref{sec:ablation-contributions}). The improvement therefore appears in descriptions of the target's research context and objectives.

The oracle input shows that the reader benefits from more informative conditioning. Supplying the target's own embedding in place of the GNN estimate improves agreement in all three content aspects, by 0.18 to 0.31, while keeping the reader unchanged (\autoref{tab:ablation-results}). This shows that the reader can use more informative conditioning to generate descriptions closer to the target paper, making improvement of the conditioning representation a promising direction.

Joint fine-tuning was intended to adapt the reader to the GNN estimates it receives in use and to let the generation loss adjust the estimator (\autoref{sec:training-joint}). It improved validation identification (\autoref{sec:results-joint}), while its final content scores remained close to those of the separately trained components, with differences within $\pm 0.03$ and no consistent direction across aspects (\autoref{tab:ablation-results}). Under the tested settings, this additional stage therefore produced no consistent further improvement in content agreement.

%----------------------------------------------------------------------------
\section{Patterns in the generated abstracts}
\label{sec:disc-patterns}
%----------------------------------------------------------------------------

The generated abstracts generally follow the conventions of scientific writing and present a coherent account of a study; 84.0\% of them receive the highest coherence score (\autoref{sec:results-final}). In stronger cases, they also preserve the target paper's research area and task, as illustrated by the land-cover classification example in \autoref{app:examples}: like the original, the generated abstract addresses the semantic segmentation of remote sensing images and reports an evaluation on the ISPRS Vaihingen dataset. Examining individual outputs shows more precisely which aspects are preserved and which change during generation.

Even in this stronger case, specific details are not supported by the original abstract. The generated text states that the model is trained on a dataset collected from Google Earth, which the original does not mention. Agreement on the task can therefore coexist with unsupported experimental details.

In other cases, a nearby task replaces the target's task. For a paper on retinal lesion segmentation and disease classification, the generated abstract remains in retinal image analysis but describes vessel segmentation. Its domain is still relevant, but the changed objective makes it a description of a different study.

The kind of contribution can also change. The original face-recognition paper analyses why recognition accuracy differs between men and women, whereas the generated abstract proposes a new face-recognition method. Proposal-style wording recurs across the outputs: the word ``novel'' occurs in 69.1\% of the generated abstracts and in 23.1\% of the original abstracts, and sixty generated abstracts begin with the same eight words, ``We present a novel approach to the problem''. This recurring wording concerns variation across different abstracts. The six-token repetition constraint instead restricts repeated token sequences within an individual output (\autoref{sec:decoding-method}).

Finally, a generated abstract can describe a recorded reference instead of the target. For a paper on exemplar normalisation, the model wrote an abstract of ResNet~\cite{he2016resnet}, which the target cites.

Recurring formulations in the training abstracts may encourage the reader to reproduce a familiar proposal-and-evaluation structure when the conditioning leaves the target's contribution ambiguous. Greedy decoding may further favour reuse of these formulations.

%----------------------------------------------------------------------------
\section{Design choices relative to the assignment}
\label{sec:disc-design}

This section explains the implementation choices summarised in \autoref{sec:scope}.

Instead of projecting the GNN estimate into the input embedding space of the language model, I map it to key and value prefixes in every layer (\autoref{sec:prefix-adapter}). Layer-wise prefixes let the paper representation condition attention directly throughout the model, rather than reaching the deeper layers only through the frozen network~\cite{liu2022ptuning}. I chose them for this reason; the experiments do not compare them with input-level prompting.

Instead of training all components together from the start, I first train the reader and the GNN estimator separately. Separate training gives each component a defined task before they are connected. The reader learns to generate abstracts from informative vectors, the target's own embedding or its reference mean (\autoref{sec:training-reader}), while the estimator learns to predict the target embedding through direct supervision (\autoref{sec:training-gnn}). Each component can therefore be trained and assessed on its own task. Joint fine-tuning then updates all trainable components together using the generation and embedding objectives (\autoref{sec:training-joint}). This joint phase is end-to-end training, and the separate stages provide its initialisation; its effect on the results is discussed in \autoref{sec:disc-representation}.

Earlier attempts to train the graph-conditioned generator jointly produced weak target identification and, in several configurations, identical abstracts for many different papers. These observations motivated training the reader and estimator separately before joint fine-tuning. The development runs differed from the final system in their training data and pipeline, so they are not a controlled comparison with it.

The reference mean provides a transparent baseline using the same reference information and reader as the GNN configuration (\autoref{sec:gnn-estimator}). Their comparison measures the benefit of learned estimation over equal-weight averaging (\autoref{sec:ablation-contributions}). Isolating the contribution of graph structure itself would require an additional learned, structure-agnostic baseline.
